\clearpage\phantomsection\label{couverture-programme}\pdfbookmark[0]{Repérage du programme}{bm-couverture-programme}\gdef\otChapNum{0}\pagestyle{otFront}
\otPageTitre{Repérage du programme}{Sixième · Programme 2025 du cycle 3}
Cette matrice relie les objectifs d’apprentissage du programme aux
exercices du livre. Les cours complètent ces repères, notamment pour
le vocabulaire, les représentations et les méthodes de construction.\par\bigskip
Les six domaines du programme sont couverts : nombres, calcul et
résolution de problèmes, y compris l’algèbre (chapitres 1 à 8) ;
grandeurs et mesures (10 à 12, et les volumes au chapitre 18) ; espace et
géométrie (13 à 18) ; organisation et gestion de données et probabilités
(19 et 20) ; proportionnalité (9, avec les échelles) ; initiation à la
pensée informatique (21, et plusieurs défis).\par\bigskip
\begin{attention}Les niveaux de difficulté et les défis sont des choix
d’entraînement. Ils ne constituent pas une liste d’obligations
supplémentaires au programme officiel.\end{attention}\clearpage
\otPageTitre{Repères et exercices}{Chapitres 1 à 3}\fontsize{10}{13}\selectfont
\otIntertitre{1. Entiers et décimaux}\begin{tabularx}{\linewidth}{@{}Yl@{}}\toprule
Grands nombres et milliards & \hyperref[exo-1-1]{1.1}, \hyperref[exo-1-2]{1.2}, \hyperref[exo-1-10]{1.10}, \hyperref[exo-1-11]{1.11}, \hyperref[exo-1-12]{1.12}\\[2mm]
Unités de numération et valeur des chiffres & \hyperref[exo-1-3]{1.3}, \hyperref[exo-1-6]{1.6}, \hyperref[exo-1-15]{1.15}\\[2mm]
Nombres décimaux et fractions décimales & \hyperref[exo-1-4]{1.4}, \hyperref[exo-1-5]{1.5}, \hyperref[exo-1-9]{1.9}, \hyperref[exo-1-13]{1.13}\\[2mm]
Écritures d’un même nombre & \hyperref[exo-1-7]{1.7}, \hyperref[exo-1-8]{1.8}, \hyperref[exo-1-14]{1.14}\\[2mm]
\bottomrule\end{tabularx}\par\vspace{4mm}
\otIntertitre{2. Repérer, comparer et arrondir}\begin{tabularx}{\linewidth}{@{}Yl@{}}\toprule
Demi-droite graduée et abscisses & \hyperref[exo-2-1]{2.1}, \hyperref[exo-2-2]{2.2}, \hyperref[exo-2-8]{2.8}\\[2mm]
Comparer et ranger & \hyperref[exo-2-3]{2.3}, \hyperref[exo-2-4]{2.4}, \hyperref[exo-2-11]{2.11}, \hyperref[exo-2-15]{2.15}\\[2mm]
Encadrer et intercaler & \hyperref[exo-2-6]{2.6}, \hyperref[exo-2-7]{2.7}, \hyperref[exo-2-13]{2.13}\\[2mm]
Arrondis et valeurs approchées & \hyperref[exo-2-5]{2.5}, \hyperref[exo-2-9]{2.9}, \hyperref[exo-2-10]{2.10}, \hyperref[exo-2-12]{2.12}, \hyperref[exo-2-14]{2.14}\\[2mm]
\bottomrule\end{tabularx}\par\vspace{4mm}
\otIntertitre{3. Calculer avec les décimaux}\begin{tabularx}{\linewidth}{@{}Yl@{}}\toprule
Additionner et soustraire des décimaux & \hyperref[exo-3-1]{3.1}, \hyperref[exo-3-2]{3.2}, \hyperref[exo-3-11]{3.11}, \hyperref[exo-3-14]{3.14}\\[2mm]
Multiplier et diviser par 10, 100, 1 000 ou 0,1… & \hyperref[exo-3-3]{3.3}, \hyperref[exo-3-4]{3.4}\\[2mm]
Produit de deux décimaux & \hyperref[exo-3-5]{3.5}, \hyperref[exo-3-6]{3.6}, \hyperref[exo-3-7]{3.7}, \hyperref[exo-3-9]{3.9}, \hyperref[exo-3-12]{3.12}\\[2mm]
Ordre de grandeur et problèmes & \hyperref[exo-3-8]{3.8}, \hyperref[exo-3-10]{3.10}, \hyperref[exo-3-13]{3.13}, \hyperref[exo-3-15]{3.15}\\[2mm]
\bottomrule\end{tabularx}\par\vspace{4mm}
\clearpage
\otPageTitre{Repères et exercices}{Chapitres 4 à 6}\fontsize{10}{13}\selectfont
\otIntertitre{4. Diviser et interpréter un partage}\begin{tabularx}{\linewidth}{@{}Yl@{}}\toprule
Sens de la division : partage ou groupement & \hyperref[exo-4-1]{4.1}, \hyperref[exo-4-13]{4.13}\\[2mm]
Division euclidienne et reste & \hyperref[exo-4-2]{4.2}, \hyperref[exo-4-3]{4.3}, \hyperref[exo-4-6]{4.6}, \hyperref[exo-4-10]{4.10}, \hyperref[exo-4-15]{4.15}\\[2mm]
Division décimale par un entier & \hyperref[exo-4-4]{4.4}, \hyperref[exo-4-5]{4.5}, \hyperref[exo-4-12]{4.12}\\[2mm]
Interpréter un quotient, un reste, une valeur approchée & \hyperref[exo-4-7]{4.7}, \hyperref[exo-4-8]{4.8}, \hyperref[exo-4-9]{4.9}, \hyperref[exo-4-11]{4.11}, \hyperref[exo-4-14]{4.14}\\[2mm]
\bottomrule\end{tabularx}\par\vspace{4mm}
\otIntertitre{5. Les fractions comme nombres}\begin{tabularx}{\linewidth}{@{}Yl@{}}\toprule
Fraction quotient et égalités à trous & \hyperref[exo-5-1]{5.1}, \hyperref[exo-5-2]{5.2}, \hyperref[exo-5-5]{5.5}, \hyperref[exo-5-13]{5.13}\\[2mm]
Fractions sur une demi-droite et graduations & \hyperref[exo-5-4]{5.4}, \hyperref[exo-5-6]{5.6}, \hyperref[exo-5-9]{5.9}\\[2mm]
Entier, décimal ou non décimal & \hyperref[exo-5-3]{5.3}, \hyperref[exo-5-7]{5.7}, \hyperref[exo-5-12]{5.12}, \hyperref[exo-5-15]{5.15}\\[2mm]
Nombres mixtes et problèmes de partage & \hyperref[exo-5-8]{5.8}, \hyperref[exo-5-10]{5.10}, \hyperref[exo-5-11]{5.11}, \hyperref[exo-5-14]{5.14}\\[2mm]
\bottomrule\end{tabularx}\par\vspace{4mm}
\otIntertitre{6. Comparer et calculer des fractions}\begin{tabularx}{\linewidth}{@{}Yl@{}}\toprule
Fractions égales et simplification & \hyperref[exo-6-1]{6.1}, \hyperref[exo-6-2]{6.2}, \hyperref[exo-6-10]{6.10}\\[2mm]
Comparer, encadrer et ranger & \hyperref[exo-6-3]{6.3}, \hyperref[exo-6-6]{6.6}, \hyperref[exo-6-9]{6.9}\\[2mm]
Additionner et soustraire & \hyperref[exo-6-4]{6.4}, \hyperref[exo-6-7]{6.7}, \hyperref[exo-6-11]{6.11}\\[2mm]
Multiplier par un entier, appliquer une fraction & \hyperref[exo-6-5]{6.5}, \hyperref[exo-6-8]{6.8}, \hyperref[exo-6-12]{6.12}, \hyperref[exo-6-13]{6.13}, \hyperref[exo-6-14]{6.14}, \hyperref[exo-6-15]{6.15}\\[2mm]
\bottomrule\end{tabularx}\par\vspace{4mm}
\clearpage
\otPageTitre{Repères et exercices}{Chapitres 7 à 9}\fontsize{10}{13}\selectfont
\otIntertitre{7. Proportions et pourcentages}\begin{tabularx}{\linewidth}{@{}Yl@{}}\toprule
Sens d’un pourcentage & \hyperref[exo-7-1]{7.1}, \hyperref[exo-7-15]{7.15}\\[2mm]
Fraction, décimal et pourcentage & \hyperref[exo-7-2]{7.2}, \hyperref[exo-7-3]{7.3}\\[2mm]
Proportion d’une partie dans un tout & \hyperref[exo-7-4]{7.4}, \hyperref[exo-7-6]{7.6}, \hyperref[exo-7-8]{7.8}, \hyperref[exo-7-9]{7.9}, \hyperref[exo-7-10]{7.10}, \hyperref[exo-7-14]{7.14}\\[2mm]
Appliquer un pourcentage & \hyperref[exo-7-5]{7.5}, \hyperref[exo-7-7]{7.7}, \hyperref[exo-7-11]{7.11}, \hyperref[exo-7-12]{7.12}, \hyperref[exo-7-13]{7.13}\\[2mm]
\bottomrule\end{tabularx}\par\vspace{4mm}
\otIntertitre{8. Chercher des inconnues et des régularités}\begin{tabularx}{\linewidth}{@{}Yl@{}}\toprule
Décrire des quantités inconnues & \hyperref[exo-8-2]{8.2}, \hyperref[exo-8-11]{8.11}\\[2mm]
Schémas en barres & \hyperref[exo-8-1]{8.1}, \hyperref[exo-8-5]{8.5}, \hyperref[exo-8-6]{8.6}, \hyperref[exo-8-7]{8.7}\\[2mm]
Raisonner sans connaître tous les nombres & \hyperref[exo-8-4]{8.4}, \hyperref[exo-8-10]{8.10}, \hyperref[exo-8-14]{8.14}\\[2mm]
Motifs évolutifs & \hyperref[exo-8-3]{8.3}, \hyperref[exo-8-8]{8.8}, \hyperref[exo-8-9]{8.9}, \hyperref[exo-8-12]{8.12}, \hyperref[exo-8-13]{8.13}, \hyperref[exo-8-15]{8.15}\\[2mm]
\bottomrule\end{tabularx}\par\vspace{4mm}
\otIntertitre{9. Proportionnalité et échelles}\begin{tabularx}{\linewidth}{@{}Yl@{}}\toprule
Reconnaître la proportionnalité & \hyperref[exo-9-1]{9.1}, \hyperref[exo-9-8]{9.8}, \hyperref[exo-9-15]{9.15}\\[2mm]
Fois plus, fois moins et sommes & \hyperref[exo-9-2]{9.2}, \hyperref[exo-9-4]{9.4}, \hyperref[exo-9-7]{9.7}, \hyperref[exo-9-10]{9.10}, \hyperref[exo-9-12]{9.12}\\[2mm]
Retour à l’unité et tableaux & \hyperref[exo-9-3]{9.3}, \hyperref[exo-9-6]{9.6}, \hyperref[exo-9-13]{9.13}\\[2mm]
Échelles & \hyperref[exo-9-5]{9.5}, \hyperref[exo-9-9]{9.9}, \hyperref[exo-9-11]{9.11}, \hyperref[exo-9-14]{9.14}\\[2mm]
\bottomrule\end{tabularx}\par\vspace{4mm}
\clearpage
\otPageTitre{Repères et exercices}{Chapitres 10 à 12}\fontsize{10}{13}\selectfont
\otIntertitre{10. Longueurs et périmètres}\begin{tabularx}{\linewidth}{@{}Yl@{}}\toprule
Unités de longueur et conversions & \hyperref[exo-10-1]{10.1}, \hyperref[exo-10-2]{10.2}, \hyperref[exo-10-5]{10.5}, \hyperref[exo-10-13]{10.13}\\[2mm]
Périmètres du carré et du rectangle & \hyperref[exo-10-3]{10.3}, \hyperref[exo-10-8]{10.8}, \hyperref[exo-10-12]{10.12}, \hyperref[exo-10-15]{10.15}\\[2mm]
Périmètre du disque & \hyperref[exo-10-4]{10.4}, \hyperref[exo-10-9]{10.9}, \hyperref[exo-10-11]{10.11}, \hyperref[exo-10-14]{10.14}\\[2mm]
Figures composées et problèmes & \hyperref[exo-10-6]{10.6}, \hyperref[exo-10-7]{10.7}, \hyperref[exo-10-10]{10.10}\\[2mm]
\bottomrule\end{tabularx}\par\vspace{4mm}
\otIntertitre{11. Comparer et mesurer des aires}\begin{tabularx}{\linewidth}{@{}Yl@{}}\toprule
Comparer des aires sans mesurer & \hyperref[exo-11-1]{11.1}, \hyperref[exo-11-7]{11.7}, \hyperref[exo-11-8]{11.8}\\[2mm]
Unités d’aire et quadrillages & \hyperref[exo-11-2]{11.2}, \hyperref[exo-11-3]{11.3}\\[2mm]
Conversions d’aires & \hyperref[exo-11-5]{11.5}, \hyperref[exo-11-6]{11.6}, \hyperref[exo-11-11]{11.11}, \hyperref[exo-11-14]{11.14}\\[2mm]
Formules et problèmes & \hyperref[exo-11-4]{11.4}, \hyperref[exo-11-9]{11.9}, \hyperref[exo-11-10]{11.10}, \hyperref[exo-11-12]{11.12}, \hyperref[exo-11-13]{11.13}, \hyperref[exo-11-15]{11.15}\\[2mm]
\bottomrule\end{tabularx}\par\vspace{4mm}
\otIntertitre{12. Horaires et durées}\begin{tabularx}{\linewidth}{@{}Yl@{}}\toprule
Lire l’heure et les unités de durée & \hyperref[exo-12-1]{12.1}, \hyperref[exo-12-2]{12.2}, \hyperref[exo-12-3]{12.3}, \hyperref[exo-12-10]{12.10}\\[2mm]
Calculer des horaires et des durées & \hyperref[exo-12-4]{12.4}, \hyperref[exo-12-5]{12.5}, \hyperref[exo-12-9]{12.9}, \hyperref[exo-12-11]{12.11}, \hyperref[exo-12-12]{12.12}\\[2mm]
Convertir des durées & \hyperref[exo-12-6]{12.6}, \hyperref[exo-12-13]{12.13}, \hyperref[exo-12-14]{12.14}\\[2mm]
Heures décimales et problèmes & \hyperref[exo-12-7]{12.7}, \hyperref[exo-12-8]{12.8}, \hyperref[exo-12-15]{12.15}\\[2mm]
\bottomrule\end{tabularx}\par\vspace{4mm}
\clearpage
\otPageTitre{Repères et exercices}{Chapitres 13 à 15}\fontsize{10}{13}\selectfont
\otIntertitre{13. Droites, distances et cercles}\begin{tabularx}{\linewidth}{@{}Yl@{}}\toprule
Vocabulaire, notations et codages & \hyperref[exo-13-1]{13.1}, \hyperref[exo-13-2]{13.2}\\[2mm]
Distance et milieu & \hyperref[exo-13-3]{13.3}, \hyperref[exo-13-8]{13.8}, \hyperref[exo-13-14]{13.14}, \hyperref[exo-13-15]{13.15}\\[2mm]
Cercle, disque, rayon, diamètre et corde & \hyperref[exo-13-4]{13.4}, \hyperref[exo-13-5]{13.5}, \hyperref[exo-13-6]{13.6}, \hyperref[exo-13-9]{13.9}\\[2mm]
Problèmes de distances & \hyperref[exo-13-7]{13.7}, \hyperref[exo-13-10]{13.10}, \hyperref[exo-13-11]{13.11}, \hyperref[exo-13-12]{13.12}, \hyperref[exo-13-13]{13.13}\\[2mm]
\bottomrule\end{tabularx}\par\vspace{4mm}
\otIntertitre{14. Médiatrice et équidistance}\begin{tabularx}{\linewidth}{@{}Yl@{}}\toprule
Définition et codage de la médiatrice & \hyperref[exo-14-1]{14.1}, \hyperref[exo-14-2]{14.2}\\[2mm]
Constructions à l’équerre et au compas & \hyperref[exo-14-3]{14.3}, \hyperref[exo-14-6]{14.6}\\[2mm]
Propriété caractéristique dans les deux sens & \hyperref[exo-14-4]{14.4}, \hyperref[exo-14-5]{14.5}, \hyperref[exo-14-7]{14.7}, \hyperref[exo-14-8]{14.8}, \hyperref[exo-14-14]{14.14}, \hyperref[exo-14-15]{14.15}\\[2mm]
Problèmes de points équidistants & \hyperref[exo-14-9]{14.9}, \hyperref[exo-14-10]{14.10}, \hyperref[exo-14-11]{14.11}, \hyperref[exo-14-12]{14.12}, \hyperref[exo-14-13]{14.13}\\[2mm]
\bottomrule\end{tabularx}\par\vspace{4mm}
\otIntertitre{15. Angles et bissectrices}\begin{tabularx}{\linewidth}{@{}Yl@{}}\toprule
Vocabulaire et types d’angles & \hyperref[exo-15-1]{15.1}, \hyperref[exo-15-2]{15.2}, \hyperref[exo-15-5]{15.5}, \hyperref[exo-15-15]{15.15}\\[2mm]
Mesurer et construire au rapporteur & \hyperref[exo-15-3]{15.3}, \hyperref[exo-15-4]{15.4}, \hyperref[exo-15-6]{15.6}, \hyperref[exo-15-12]{15.12}\\[2mm]
Angles adjacents, opposés et supplémentaires & \hyperref[exo-15-7]{15.7}, \hyperref[exo-15-8]{15.8}, \hyperref[exo-15-9]{15.9}, \hyperref[exo-15-11]{15.11}\\[2mm]
Bissectrice et problèmes & \hyperref[exo-15-10]{15.10}, \hyperref[exo-15-13]{15.13}, \hyperref[exo-15-14]{15.14}\\[2mm]
\bottomrule\end{tabularx}\par\vspace{4mm}
\clearpage
\otPageTitre{Repères et exercices}{Chapitres 16 à 18}\fontsize{10}{13}\selectfont
\otIntertitre{16. Triangles et cercle circonscrit}\begin{tabularx}{\linewidth}{@{}Yl@{}}\toprule
Construire des triangles & \hyperref[exo-16-5]{16.5}, \hyperref[exo-16-6]{16.6}, \hyperref[exo-16-7]{16.7}, \hyperref[exo-16-8]{16.8}\\[2mm]
Triangles particuliers & \hyperref[exo-16-2]{16.2}, \hyperref[exo-16-3]{16.3}, \hyperref[exo-16-4]{16.4}, \hyperref[exo-16-9]{16.9}\\[2mm]
Somme des angles & \hyperref[exo-16-1]{16.1}, \hyperref[exo-16-13]{16.13}, \hyperref[exo-16-14]{16.14}, \hyperref[exo-16-15]{16.15}\\[2mm]
Médiatrices et cercle circonscrit & \hyperref[exo-16-10]{16.10}, \hyperref[exo-16-11]{16.11}, \hyperref[exo-16-12]{16.12}\\[2mm]
\bottomrule\end{tabularx}\par\vspace{4mm}
\otIntertitre{17. Symétrie axiale}\begin{tabularx}{\linewidth}{@{}Yl@{}}\toprule
Axes de symétrie d’une figure & \hyperref[exo-17-1]{17.1}, \hyperref[exo-17-14]{17.14}, \hyperref[exo-17-15]{17.15}\\[2mm]
Constructions sur quadrillage et papier pointé & \hyperref[exo-17-2]{17.2}, \hyperref[exo-17-3]{17.3}, \hyperref[exo-17-4]{17.4}, \hyperref[exo-17-5]{17.5}, \hyperref[exo-17-9]{17.9}\\[2mm]
Constructions sur papier uni & \hyperref[exo-17-6]{17.6}, \hyperref[exo-17-7]{17.7}, \hyperref[exo-17-8]{17.8}, \hyperref[exo-17-10]{17.10}, \hyperref[exo-17-13]{17.13}\\[2mm]
Propriétés de la symétrie & \hyperref[exo-17-11]{17.11}, \hyperref[exo-17-12]{17.12}\\[2mm]
\bottomrule\end{tabularx}\par\vspace{4mm}
\otIntertitre{18. Solides, assemblages et volumes}\begin{tabularx}{\linewidth}{@{}Yl@{}}\toprule
Reconnaître les solides usuels & \hyperref[exo-18-1]{18.1}, \hyperref[exo-18-2]{18.2}\\[2mm]
Vues et perspectives & \hyperref[exo-18-4]{18.4}, \hyperref[exo-18-7]{18.7}, \hyperref[exo-18-14]{18.14}\\[2mm]
Dénombrer des assemblages de cubes & \hyperref[exo-18-3]{18.3}, \hyperref[exo-18-5]{18.5}, \hyperref[exo-18-6]{18.6}, \hyperref[exo-18-8]{18.8}, \hyperref[exo-18-10]{18.10}, \hyperref[exo-18-12]{18.12}, \hyperref[exo-18-13]{18.13}\\[2mm]
Volumes en centimètres cubes & \hyperref[exo-18-9]{18.9}, \hyperref[exo-18-11]{18.11}, \hyperref[exo-18-15]{18.15}\\[2mm]
\bottomrule\end{tabularx}\par\vspace{4mm}
\clearpage
\otPageTitre{Repères et exercices}{Chapitres 19 à 21}\fontsize{10}{13}\selectfont
\otIntertitre{19. Données et enquêtes}\begin{tabularx}{\linewidth}{@{}Yl@{}}\toprule
Planifier une enquête et recueillir des données & \hyperref[exo-19-1]{19.1}, \hyperref[exo-19-14]{19.14}\\[2mm]
Construire un tableau de données & \hyperref[exo-19-2]{19.2}, \hyperref[exo-19-6]{19.6}, \hyperref[exo-19-7]{19.7}, \hyperref[exo-19-10]{19.10}\\[2mm]
Lire tableaux, diagrammes et courbes & \hyperref[exo-19-3]{19.3}, \hyperref[exo-19-4]{19.4}, \hyperref[exo-19-5]{19.5}, \hyperref[exo-19-13]{19.13}\\[2mm]
Filtrer et interpréter avec esprit critique & \hyperref[exo-19-8]{19.8}, \hyperref[exo-19-9]{19.9}, \hyperref[exo-19-11]{19.11}, \hyperref[exo-19-12]{19.12}, \hyperref[exo-19-15]{19.15}\\[2mm]
\bottomrule\end{tabularx}\par\vspace{4mm}
\otIntertitre{20. Comprendre le hasard}\begin{tabularx}{\linewidth}{@{}Yl@{}}\toprule
Probabilité entre 0 et 1 & \hyperref[exo-20-1]{20.1}, \hyperref[exo-20-5]{20.5}, \hyperref[exo-20-10]{20.10}\\[2mm]
Calculer en situation d’équiprobabilité & \hyperref[exo-20-2]{20.2}, \hyperref[exo-20-4]{20.4}, \hyperref[exo-20-6]{20.6}, \hyperref[exo-20-7]{20.7}, \hyperref[exo-20-11]{20.11}\\[2mm]
Écritures et comparaisons & \hyperref[exo-20-3]{20.3}, \hyperref[exo-20-9]{20.9}, \hyperref[exo-20-12]{20.12}, \hyperref[exo-20-14]{20.14}, \hyperref[exo-20-15]{20.15}\\[2mm]
Fréquences et essais répétés & \hyperref[exo-20-8]{20.8}, \hyperref[exo-20-13]{20.13}\\[2mm]
\bottomrule\end{tabularx}\par\vspace{4mm}
\otIntertitre{21. Instructions, répétitions et chemins}\begin{tabularx}{\linewidth}{@{}Yl@{}}\toprule
Instruction, séquence, entrée et sortie & \hyperref[exo-21-1]{21.1}, \hyperref[exo-21-3]{21.3}\\[2mm]
Exécuter un programme à la main & \hyperref[exo-21-2]{21.2}, \hyperref[exo-21-5]{21.5}, \hyperref[exo-21-9]{21.9}\\[2mm]
Répétitions & \hyperref[exo-21-6]{21.6}, \hyperref[exo-21-7]{21.7}, \hyperref[exo-21-10]{21.10}, \hyperref[exo-21-12]{21.12}, \hyperref[exo-21-13]{21.13}, \hyperref[exo-21-14]{21.14}\\[2mm]
Coder un chemin & \hyperref[exo-21-4]{21.4}, \hyperref[exo-21-8]{21.8}, \hyperref[exo-21-11]{21.11}, \hyperref[exo-21-15]{21.15}\\[2mm]
\bottomrule\end{tabularx}\par\vspace{4mm}
\clearpage
